% answers_short.tex —— 2009 年·普通高等学校招生考试·新课标卷（琼、宁卷）·文科 · 简短答案速查
% 由 tools/make_short_answers.py 自 answers.tex 的 \gksolhead/\gkanswer 提取，
% 并与 2009/2009_数学真题_新课标文_答案要点.md 逐题比对后生成，勿手改（改源 MD 后重跑）。

\begin{gknotice}
\textbf{使用说明}：本节只给最终答案便于快速核对，完整推导见解析版。答案为编写者独立推导，\textbf{非官方标准答案}。
\end{gknotice}

\gkpart{选择题}
\begin{center}\begin{tabular}{ccccccccccccc}
\toprule
题号 & 1 & 2 & 3 & 4 & 5 & 6 & 7 & 8 & 9 & 10 & 11 & 12 \\
\midrule
答案 & D & C & C & A & B & B & A & C & D & B & A & C \\
\bottomrule
\end{tabular}\end{center}

\gkpart{填空题}
\begin{itemize}[leftmargin=2.6em]
  \item[\textbf{13.}] $y=3x+1$。
  \item[\textbf{14.}] $y^2=4x$。
  \item[\textbf{15.}] $\frac{15}{2}$。
  \item[\textbf{16.}] $0$。
\end{itemize}

\gkpart{解答题}
\begin{itemize}[leftmargin=2.6em]
  \item[\textbf{17.}] $\frac{16}{65}$。
  \item[\textbf{18.}] （1）证明见解析；（2）$\frac83$。
  \item[\textbf{19.}] （1）$\frac1{100}$；（2）$x=5$，$y=15$，$B$ 类工人中个体间的差异程度更小；$A$ 类、$B$ 类以及全厂工人生产能力的平均数估计值分别为 $123$、$133.8$、$131.1$。
  \item[\textbf{20.}] （1）$\frac{x^2}{16}+\frac{y^2}{7}=1$；（2）轨迹方程为 $y=\pm\frac{4\sqrt7}{3}$（$-4\leqslant x\leqslant4$），轨迹是两条平行于 $x$ 轴的线段。
  \item[\textbf{21.}] （1）极大值为 $f(-1)=6$，极小值为 $f(3)=-26$；（2）$\frac14<a\leqslant\frac45$。
\end{itemize}

\gkpart{选做题}
\begin{itemize}[leftmargin=2.6em]
  \item[\textbf{22.}] （1）（2）证明见解析。
  \item[\textbf{23.}] （1）$C_1:(x+4)^2+(y-3)^2=1$，表示圆；$C_2:\frac{x^2}{64}+\frac{y^2}{9}=1$，表示椭圆；（2）$\frac{8\sqrt5}{5}$。
  \item[\textbf{24.}] （1）$y=\begin{cases}160-10x,&0\leqslant x\leqslant10,\\ 80-2x,&10\leqslant x\leqslant20,\\ 10x-160,&20\leqslant x\leqslant30;\end{cases}$（2）$9\leqslant x\leqslant23$。
\end{itemize}
